% ============================================================
% ABSTRACT
% ============================================================
\begin{abstract}
Large language model (LLM) role-playing agents face a fundamental tension: maintaining consistent character identity while exhibiting emotionally diverse responses across multi-turn conversations. Existing approaches either enforce rigid persona descriptions (producing emotionally flat output) or employ self-corrective reflection (which degrades consistency). We propose the \textbf{Thermodynamic Persona Engine (TPE)}, a physics-inspired framework that couples a frustration-driven ``temperature'' to behavioral signal noise, while using mass-weighted memory retrieval with temporal decay to anchor response style. In a controlled evaluation across three LLMs (225 experiments, 30 turns each, 5 scenarios), TPE achieves a 32\% improvement in emotional variance on qwen3-max (Bonferroni-adjusted $p = 0.008$, Cliff's $d = +0.75$, large effect) without reducing personality consistency. We additionally report two generalizable findings: (1) Reflexion-style self-correction significantly degrades persona consistency across all models ($p < 0.001$ on 2/3 models), and (2) strong RLHF alignment creates a hard ceiling on prompt-level persona control---on gpt-5-mini, anti-alignment instructions trigger safety rollbacks that amplify template counselor behavior (Collapse Index $22\times$ higher than baseline). Code and data are available at [anonymized].
\end{abstract}
